\documentclass[10pt,a4paper]{amsart}
\usepackage[margin=19mm]{geometry}
\usepackage{amsmath,amssymb,mathtools}
\usepackage[scr=rsfs,cal=euler]{mathalfa}
\usepackage{xfrac}
\usepackage[unicode,hidelinks,bookmarks=false]{hyperref}
\newcommand{\ie}{i.e.\ }
\newcommand{\cf}{cf.\ }
\newcommand{\resp}{respectively\ }
\newcommand{\Subsection}{\subsection}
\providecommand{\nobreakdash}{\nobreak\hbox{-}}
\setlength{\emergencystretch}{2em}
\multlinegap0pt
\usepackage{xcolor}
\usepackage{tikz}
\usepackage{amssymb}
\newtheorem{lemma}{Lemma}
\title{Raimi's theorem for manifolds with circle symmetry}
\author{Dung The Tran}
\date{Source: arXiv:2512.08676v2 (CC BY 4.0)}
\begin{document}
\maketitle

\noindent\emph{Source and licence.} Raimi's theorem for manifolds with circle symmetry. Version arXiv:2512.08676v2 (CC BY 4.0), distributed by arXiv under the Creative Commons Attribution 4.0 International licence, http://creativecommons.org/licenses/by/4.0/, as recorded in the arXiv licence metadata for this exact version and shown on the abstract page https://arxiv.org/abs/2512.08676v2. Full licence notice: \texttt{licenses/arxiv-2512.08676-CC-BY-4.0.txt}.\par\smallskip

\noindent\textit{Editorial scope.} Counted unit: the article's lemma lem:Haar-fibres (source0.tex lines 622--630). The article's complete original proof of it is delivered with it (source0.tex lines 631--637, one \texttt{proof} environment of 7 lines). No mathematics is written, repaired or re-derived: the statement and the proof are the article's own lines, in the article's own notation, and no retained line is substituted or re-flowed.  No further source-local context is needed: every cross-reference of every retained line resolves inside the two delivered spans.  Nothing was omitted from any retained span, so the package carries no omission record.  No retained line contains a citation, so the wrapper carries no bibliography and no bibliography entry.  No retained line contains a figure environment, an image-inclusion command or drawing code, so no image is delivered and the package declares no asset dependency.  Editorial formatting only, in addition: an AMS \texttt{amsart} preamble that restates the article's own declarations and macros used by the retained lines, namely the environments \texttt{lemma} on separate counters, together with the standard packages those lines need.  The retained environments print in this document as Lemma 1 in order of appearance, whereas the article prints the same items with the numbering of its own full document.  The complete native source of the article as distributed in the arXiv e-print for this exact version is bundled unchanged as \texttt{sources/190679/source0.tex}.
\begin{lemma}[Invariant conditional measures on circle fibres]\label{lem:Haar-fibres}
Let $(X,\mu)$ be a probability space equipped with a measurable, measure-preserving action of the circle group $\mathbb{T}$.
Suppose that $\pi:X\to Y$ is a measurable factor map whose fibres $\pi^{-1}(\{y\})$ are $\mathbb{T}$--orbits, and let $\{\mu_y\}_{y\in Y}$ be a Rokhlin disintegration of $\mu$ over $Y$.
Then for $\nu$--almost every $y\in Y$, the conditional measure $\mu_y$ is invariant under translations on the fibre. Consequently,
\[
\mu_y = \mu_1 \qquad \text{for $\nu$--a.e.\ } y\in Y,
\]
where $\mu_1$ denotes the normalized Haar measure on $\mathbb{T}$.
\end{lemma}

\begin{proof}
Since $\mu$ is invariant under the $\mathbb{T}$--action, the Rokhlin
conditional measures $\mu_y$ must be invariant under translations
along the fibres for $\nu$--almost every $y$.
By uniqueness of the translation-invariant probability measure on
$\mathbb{T}$, it follows that $\mu_y=\mu_1$.
\end{proof}

\end{document}
